\unit{Damping, Langevin noise and master equations}
\unitref{W5L2; 2024 Q4; 2025 Q4}
\subsection{Two ways of describing the same physics}
The Heisenberg picture evolves operators;
the Schr\"odinger picture evolves density matrices.
For a closed system,
\[
\dot O=\frac{i}{\hbar}[H,O],\qquad
\dot\rho=-\frac{i}{\hbar}[H,\rho].
\]
An open system also has damping and fluctuations,
because it couples to many unobserved bath modes.
The Markov approximation replaces bath memory
by effectively instantaneous noise; it is
useful when bath correlations decay quickly
compared with the system dynamics.

\subsection{Where the quantum Langevin equation comes from}
A cavity mode \(a\) coupled weakly to
a continuum of oscillators \(b_\omega\)
has a Hamiltonian of the schematic form
\[
H/\hbar=\omega_ca^\dagger a+
\int\omega b_\omega^\dagger b_\omega d\omega+
i\int\kappa(\omega)
(a^\dagger b_\omega-b_\omega^\dagger a)d\omega.
\]
One first solves the bath equation for \(b_\omega(t)\);
substituting it back into the cavity equation
produces a \emph{damping} term and an
\emph{input noise} term. In a frame rotating
at the cavity resonance,
\[
\boxed{\dot a=-\frac\gamma2a+
\sqrt\gamma\,a_{\rm in}(t)}
\]
for an undriven cavity with one damping port.
The noise operator is not an optional
classical disturbance: it ensures that
operator commutators remain physically valid.

\subsection{Fourier solving is usually easiest}
Define
\(\tilde a(\omega)=\int_{-\infty}^\infty
a(t)e^{-i\omega t}\,dt\), so
\(\widetilde{\dot a}=i\omega\tilde a\)
when boundary terms vanish.
Then
\[
(i\omega+\gamma/2)\tilde a
=\sqrt\gamma\,\tilde a_{\rm in},
\qquad
\boxed{\tilde a=
\frac{\sqrt\gamma}{\gamma/2+i\omega}
\tilde a_{\rm in}}.
\]
The denominator is the cavity susceptibility.
Near zero detuning it has a large response;
far from resonance the response is smaller.
A consistent Fourier convention is essential
because another sign convention replaces
\(i\omega\) by \(-i\omega\).

\subsection{Input--output theory}
With one common convention,
\[
a_{\rm out}=\sqrt\gamma\,a-a_{\rm in}.
\]
Eliminating the intracavity operator gives
a frequency-dependent relation between
input and output fields. For a two-port
cavity with \(\gamma=\gamma_1+\gamma_2\),
\[
\dot a=-\frac\gamma2a+
\sqrt{\gamma_1}a_{{\rm in},1}
+\sqrt{\gamma_2}a_{{\rm in},2},
\]
\[
\boxed{
\tilde a_{{\rm out},1}
=\frac{\gamma_1-\gamma/2-i\omega}
{\gamma/2+i\omega}\tilde a_{{\rm in},1}
+\frac{\sqrt{\gamma_1\gamma_2}}
{\gamma/2+i\omega}\tilde a_{{\rm in},2}.}
\]
At \(\omega=0\) and \(\gamma_1=\gamma_2\)
the coherent reflection coefficient vanishes.
But the second input's vacuum fluctuations
must remain in the \emph{field operator}.

\subsection{The Lindblad master equation}
After tracing over a Markovian bath, a general
trace-preserving reduced evolution is
\[
\boxed{\dot\rho=-\frac{i}{\hbar}[H,\rho]+
\sum_j\gamma_j\mathcal D[L_j]\rho},
\]
where
\[
\mathcal D[L]\rho=L\rho L^\dagger
-\frac12(L^\dagger L\rho+\rho L^\dagger L).
\]
A cavity losing photons into vacuum has
\(L=a\), while a two-level atom decaying
spontaneously has \(L=\sigma_-\).
Pure dephasing can be modelled with a
Hermitian jump operator related to \(\sigma_z\).
The jump operator tells you \emph{which
physical transition the environment causes}.

\begin{worked}{Photon number decays exponentially}
For zero-temperature cavity loss,
\(\dot\rho=\gamma\mathcal D[a]\rho\).
Use \(\frac{d}{dt}\langle n\rangle
=\operatorname{Tr}(n\dot\rho)\) and
\([n,a]=-a\). After rearrangement,
\[
\frac{d}{dt}\langle n\rangle
=-\gamma\langle n\rangle.
\]
Hence
\[
\boxed{\langle n(t)\rangle
=\langle n(0)\rangle e^{-\gamma t}}.
\]
The field amplitude decays at half
that exponent (\(\gamma/2\)), explaining
the factor 2 in the Langevin equation.
\end{worked}

\subsection{Parametric squeezing in a cavity}
For interaction
\(H_I=(i\hbar\chi/2)(a^2-a^{\dagger2})\),
the Heisenberg commutator contributes
\(-\chi a^\dagger\) to \(\dot a\).
Define \(X_+=a+a^\dagger\), giving
\[
\dot X_+=-(\gamma/2+\chi)X_+
+\sqrt\gamma X_{+,\rm in}.
\]
The output noise transfer factor leads to
\[
\boxed{V_+(\omega)=
\frac{(\gamma/2-\chi)^2+\omega^2}
{(\gamma/2+\chi)^2+\omega^2}}
\]
for unit vacuum input spectrum and
\(0<\chi<\gamma/2\).
The \emph{variance} comes from
field-fluctuation correlations, not
the square of the mean field.

\begin{checkpoint}
For a Langevin exam question:
identify the Hamiltonian and port rates,
calculate the commutator, Fourier transform,
apply the exact input--output sign
convention, then evaluate the noise spectrum.
\end{checkpoint}
